\par\Needspace{4\baselineskip}
\begin{center}
\hypertarget{native-067-article050}{}\pdfbookmark[2]{50.}{article-050}50.
\end{center}
Der Ort eines beliebigen Punktes auf der Oberfläche der Himmelskugel lässt sich in Bezug auf die Ekliptik am einfachsten durch seinen Abstand von der Ekliptik (seine \emph{Breite}) und durch den Abstand des Lotfußpunkts vom Äquinoktialpunkt (seine \emph{Länge}) bestimmen. Der Lotfußpunkt liegt dort, wo das vom gegebenen Punkt gefällte Lot die Ekliptik trifft. Die Breite wird auf beiden Seiten der Ekliptik bis
\SourceLeaf{068}$90^0$ gemessen und im nördlichen Bereich als positiv, im südlichen als negativ betrachtet. Dem heliozentrischen Ort des Himmelskörpers, das heißt der Projektion der von der Sonne zum Körper gezogenen Geraden auf die Himmelskugel, sollen die Länge $\lambda$ und die Breite $\beta$ entsprechen. Ferner sei $u$ der Abstand des heliozentrischen Ortes vom aufsteigenden Knoten (das \emph{Breitenargument}), $i$ die Neigung der Bahn und $\AscendingNode$ die Länge des aufsteigenden Knotens. Dann sind $i$, $u$, $\beta$ und $\lambda-\AscendingNode$ Bestandteile eines rechtwinkligen sphärischen Dreiecks, zwischen denen die folgenden Beziehungen bestehen. Wie sich leicht zeigen lässt, gelten sie ohne jede Einschränkung.
\[
\begin{array}{rl}
\mathrm{I.}&\operatorname{tang}(\lambda-\AscendingNode)=\cos i\operatorname{tang}u\\[4pt]
\mathrm{II.}&\operatorname{tang}\beta=\operatorname{tang}i\sin(\lambda-\AscendingNode)\\[4pt]
\mathrm{III.}&\sin\beta=\sin i\sin u\\[4pt]
\mathrm{IV.}&\cos u=\cos\beta\cos(\lambda-\AscendingNode).
\end{array}
\]
Sind $i$ und $u$ gegeben, so lässt sich daraus $\lambda-\AscendingNode$ nach Gleichung I bestimmen und dann $\beta$ nach II oder III, sofern $\beta$ nicht zu nahe an $\pm90^0$ liegt. Formel IV kann bei Bedarf zur Kontrolle der Rechnung verwendet werden. Die Formeln I und IV zeigen außerdem, dass $\lambda-\AscendingNode$ und $u$ stets im selben Quadranten liegen, wenn $i$ zwischen $0$ und $90^0$ liegt. Dagegen liegen $\lambda-\AscendingNode$ und $360^0-u$ immer dann im selben Quadranten, wenn $i$ zwischen $90^0$ und $180^0$ liegt, also nach üblichem Sprachgebrauch immer dann, wenn die Bewegung rückläufig ist. Damit löst sich die Mehrdeutigkeit, die bei der Bestimmung von $\lambda-\AscendingNode$ über den Tangens nach Formel I verbleibt, von selbst auf.

Die folgenden Formeln ergeben sich leicht durch Kombination der vorangehenden.
\[
\begin{array}{rl}
\mathrm{V.}&\sin(u-\lambda+\AscendingNode)=2\sin\tfrac{1}{2}i^2\sin u\cos(\lambda-\AscendingNode)\\[4pt]
\mathrm{VI.}&\sin(u-\lambda+\AscendingNode)=\operatorname{tang}\tfrac{1}{2}i\sin\beta\cos(\lambda-\AscendingNode)\\[4pt]
\mathrm{VII.}&\sin(u-\lambda+\AscendingNode)=\operatorname{tang}\tfrac{1}{2}i\operatorname{tang}\beta\cos u\\[4pt]
\mathrm{VIII.}&\sin(u+\lambda-\AscendingNode)=2\cos\tfrac{1}{2}i^2\sin u\cos(\lambda-\AscendingNode)\\[4pt]
\mathrm{IX.}&\sin(u+\lambda-\AscendingNode)=\operatorname{cotang}\tfrac{1}{2}i\sin\beta\cos(\lambda-\AscendingNode)\\[4pt]
\mathrm{X.}&\sin(u+\lambda-\AscendingNode)=\operatorname{cotang}\tfrac{1}{2}i\operatorname{tang}\beta\cos u.
\end{array}
\]
Der Winkel $u-\lambda+\AscendingNode$, wenn $i$ unter $90^0$ liegt, oder $u+\lambda-\AscendingNode$, wenn $i$ über $90^0$ liegt, heißt nach üblichem Sprachgebrauch \emph{Reduktion auf die Ekliptik}. Er ist nämlich die Differenz zwischen der heliozentrischen Länge $\lambda$ und der Länge in der Bahn, die nach diesem Sprachgebrauch $\AscendingNode\pm u$ beträgt (nach unserem $\AscendingNode+u$). Ist
\SourceLeaf{069}die Neigung entweder klein oder nur wenig von $180^0$ verschieden, so kann man diese Reduktion als Größe zweiter Ordnung auffassen. Gerade in diesem Fall ist es vorzuziehen, zunächst $\beta$ nach Formel III und dann $\lambda$ nach VII oder X zu berechnen. So lässt sich eine höhere Genauigkeit erzielen als mit Formel I.

Fällt man vom Ort des Himmelskörpers im Raum ein Lot auf die Ebene der Ekliptik, so heißt der Abstand des Schnittpunkts von der Sonne die \emph{verkürzte Entfernung}. Bezeichnet man diese mit $r'$ und den Radiusvektor mit $r$, so erhält man
\[
\begin{array}{rl}
\mathrm{XI.}&r'=r\cos\beta.
\end{array}
\]
